% GENERATED FILE. Edit canonical lessons, then run npm run generate:books.
% canonical-source-hash: fnv1a64:8c17b9d22eee8320
% canonical-lessons: SA-C501-grhasthah, SA-C501-sambandhah, SA-C501-sangacchate, SA-C501-dampati, SA-C501-pitarau, SA-R501-review-495-501, SA-R501-review-final-501

\chapter{Getting Along}
\label{ch:sa-getting-along}
\hlchaptermodality{501}

\begin{chapteropening}
\textbf{By the end of this chapter:} \emph{I can say how family members get on.}

Complete the last lesson of chapter 501: I can say how family members get on.
\end{chapteropening}

\section[\texorpdfstring{\sk{गृहस्थः} (gṛhasthaḥ)}{गृहस्थः (gṛhasthaḥ)}]{\sk{गृहस्थः} (gṛhasthaḥ) — a householder}

\label{lesson:SA-C501-grhasthah}



\begin{quote}
Before the new one: say the Sanskrit for a great-grandfather, then the Sanskrit for offspring.
\end{quote}

\subsection*{You'll want to know: \sk{गृहस्थः}}
\textbf{\sk{गृहस्थः}} — \emph{gṛhasthaḥ} — \textquotedblleft{}a householder\textquotedblright{}.

\begin{quote}
\sk{गृहस्थः} \sk{कुटुम्बं} \sk{पालयति।}
\end{quote}

\begin{quote}
\emph{gṛhasthaḥ kuṭumbaṃ pālayati.}
\end{quote}

\begin{quote}
The householder takes care of the family.
\end{quote}

A masculine noun.

\begin{cousinweb}
\textbf{\sk{गृहस्थः}} (\emph{gṛhasthaḥ}) is built from \textbf{\sk{गृहम्}}, \textquotedblleft{}a house\textquotedblright{}, and \emph{stha}, \textquotedblleft{}staying\textquotedblright{}, from the root √sthā of \textbf{\sk{तिष्ठति}}; MW gives \textquotedblleft{}a householder\textquotedblright{}, a Brāhman in the second stage of his life (MW).
\end{cousinweb}

\begin{grammarlens}[title={What you've built}]
One more word for describing what you lived and felt.
\end{grammarlens}

\subsection*{Your turn}
\begin{itemize}
\raggedright
  \item \emph{Say it:} \emph{gṛhasthaḥ}
  \item \emph{Say it:} \emph{gṛhasthaḥ}, once more
  \item \emph{Say it:} \emph{gṛhasthaḥ kuṭumbaṃ pālayati.} — The householder takes care of the family.
  \item \emph{Recall:} say the Sanskrit for a great-grandfather, then the Sanskrit for offspring, then say \emph{gṛhasthaḥ} again
\end{itemize}

\begin{tcolorbox}[breakable,colback=teal!4,colframe=teal!35!black,title={Before you move on}]
What does \sk{गृहस्थः} mean? (\textquotedblleft{}A householder\textquotedblright{}.) Say it once more.
\end{tcolorbox}
\section[\texorpdfstring{\sk{सम्बन्धः} (sambandhaḥ)}{सम्बन्धः (sambandhaḥ)}]{\sk{सम्बन्धः} (sambandhaḥ) — a relationship}

\label{lesson:SA-C501-sambandhah}



\begin{quote}
Before the new one: say the Sanskrit for offspring, then the Sanskrit for a householder.
\end{quote}

\subsection*{You'll want to know: \sk{सम्बन्धः}}
\textbf{\sk{सम्बन्धः}} — \emph{sambandhaḥ} — \textquotedblleft{}a relationship\textquotedblright{}.

\begin{quote}
\sk{मम} \sk{मातुलेन} \sk{सह} \sk{मम} \sk{सम्बन्धः} \sk{मधुरः} \sk{अस्ति।}
\end{quote}

\begin{quote}
\emph{mama mātulena saha mama sambandhaḥ madhuraḥ asti.}
\end{quote}

\begin{quote}
I have a sweet relationship with my uncle.
\end{quote}

A masculine noun.

\begin{cousinweb}
\textbf{\sk{सम्बन्धः}} (\emph{sambandhaḥ}) is built from \emph{sam} and the root √bandh, \textquotedblleft{}to bind\textquotedblright{}, the root of \textbf{\sk{बध्नाति}}; MW gives \textquotedblleft{}close connection, relation\textquotedblright{} and \textquotedblleft{}a relative\textquotedblright{} (MW).
\end{cousinweb}

\begin{grammarlens}[title={What you've built}]
One more word for describing what you lived and felt.
\end{grammarlens}

\subsection*{Your turn}
\begin{itemize}
\raggedright
  \item \emph{Say it:} \emph{sambandhaḥ}
  \item \emph{Say it:} \emph{sambandhaḥ}, once more
  \item \emph{Say it:} \emph{mama mātulena saha mama sambandhaḥ madhuraḥ asti.} — I have a sweet relationship with my uncle.
  \item \emph{Recall:} say the Sanskrit for offspring, then the Sanskrit for a householder, then say \emph{sambandhaḥ} again
\end{itemize}

\begin{tcolorbox}[breakable,colback=teal!4,colframe=teal!35!black,title={Before you move on}]
What does \sk{सम्बन्धः} mean? (\textquotedblleft{}A relationship\textquotedblright{}.) Say it once more.
\end{tcolorbox}
\section[\texorpdfstring{\sk{संगच्छते} (saṅgacchate)}{संगच्छते (saṅgacchate)}]{\sk{संगच्छते} (saṅgacchate) — gets along}

\label{lesson:SA-C501-sangacchate}



\begin{quote}
Before the new one: say the Sanskrit for a householder, then the Sanskrit for a relationship.
\end{quote}

\subsection*{You'll want to know: \sk{संगच्छते}}
\textbf{\sk{संगच्छते}} — \emph{saṅgacchate} — \textquotedblleft{}gets along\textquotedblright{}.

\begin{quote}
\sk{अरुणः} \sk{भ्रात्रा} \sk{सह} \sk{सम्यक्} \sk{संगच्छते।}
\end{quote}

\begin{quote}
\emph{aruṇaḥ bhrātrā saha samyak saṅgacchate.}
\end{quote}

\begin{quote}
Arun gets along well with his brother.
\end{quote}

Past (he …): \sk{सः} \sk{भ्रात्रा} \sk{सह} \sk{संगतवान्}, \emph{saḥ bhrātrā saha saṅgatavān}.

\begin{cousinweb}
\textbf{\sk{संगच्छते}} (\emph{saṅgacchate}) is built from \emph{sam}, \textquotedblleft{}together\textquotedblright{}, and the root √gam, \textquotedblleft{}to go\textquotedblright{}, the root of \textbf{\sk{गच्छति}}: to go together; MW's ppp \emph{saṃgata} is \textquotedblleft{}met, united, allied with\textquotedblright{} (MW).
\end{cousinweb}

\begin{grammarlens}[title={What you've built}]
One more word for describing what you lived and felt.
\end{grammarlens}

\subsection*{Your turn}
\begin{itemize}
\raggedright
  \item \emph{Say it:} \emph{saṅgacchate}
  \item \emph{Say it:} \emph{saṅgacchate}, once more
  \item \emph{Say it:} \emph{aruṇaḥ bhrātrā saha samyak saṅgacchate.} — Arun gets along well with his brother.
  \item \emph{Recall:} say the Sanskrit for a householder, then the Sanskrit for a relationship, then say \emph{saṅgacchate} again
\end{itemize}

\begin{tcolorbox}[breakable,colback=teal!4,colframe=teal!35!black,title={Before you move on}]
What does \sk{संगच्छते} mean? (\textquotedblleft{}Gets along\textquotedblright{}.) Say it once more.
\end{tcolorbox}
\section[\texorpdfstring{\sk{दम्पती} (dampatī)}{दम्पती (dampatī)}]{\sk{दम्पती} (dampatī) — husband and wife}

\label{lesson:SA-C501-dampati}



\begin{quote}
Before the new one: say the Sanskrit for a relationship, then the Sanskrit for gets along.
\end{quote}

\subsection*{You'll want to know: \sk{दम्पती}}
\textbf{\sk{दम्पती}} — \emph{dampatī} — \textquotedblleft{}husband and wife\textquotedblright{}.

\begin{quote}
\sk{दम्पती} \sk{प्रतिदिनं} \sk{सह} \sk{खादतः।}
\end{quote}

\begin{quote}
\emph{dampatī pratidinaṃ saha khādataḥ.}
\end{quote}

\begin{quote}
The couple eat together every day.
\end{quote}

A masculine noun, used in the dual.

\begin{cousinweb}
\textbf{\sk{दम्पती}} (\emph{dampatī}) is built from \emph{dam}, \textquotedblleft{}house\textquotedblright{}, and \textbf{\sk{पतिः}}, \textquotedblleft{}master\textquotedblright{}; MW gives the dual \textquotedblleft{}the two masters, husband and wife\textquotedblright{} (MW).
\end{cousinweb}

\begin{grammarlens}[title={What you've built}]
One more word for describing what you lived and felt.
\end{grammarlens}

\subsection*{Your turn}
\begin{itemize}
\raggedright
  \item \emph{Say it:} \emph{dampatī}
  \item \emph{Say it:} \emph{dampatī}, once more
  \item \emph{Say it:} \emph{dampatī pratidinaṃ saha khādataḥ.} — The couple eat together every day.
  \item \emph{Recall:} say the Sanskrit for a relationship, then the Sanskrit for gets along, then say \emph{dampatī} again
\end{itemize}

\begin{tcolorbox}[breakable,colback=teal!4,colframe=teal!35!black,title={Before you move on}]
What does \sk{दम्पती} mean? (\textquotedblleft{}Husband and wife\textquotedblright{}.) Say it once more.
\end{tcolorbox}
\section[\texorpdfstring{\sk{पितरौ} (pitarau)}{पितरौ (pitarau)}]{\sk{पितरौ} (pitarau) — parents}

\label{lesson:SA-C501-pitarau}



\begin{quote}
Before the new one: say the Sanskrit for gets along, then the Sanskrit for husband and wife.
\end{quote}

\subsection*{You'll want to know: \sk{पितरौ}}
\textbf{\sk{पितरौ}} — \emph{pitarau} — \textquotedblleft{}parents\textquotedblright{}.

\begin{quote}
\sk{अहं} \sk{प्रतिदिनं} \sk{पितरौ} \sk{नमामि।}
\end{quote}

\begin{quote}
\emph{ahaṃ pratidinaṃ pitarau namāmi.}
\end{quote}

\begin{quote}
Every day I bow to my parents.
\end{quote}

A masculine noun, used in the dual.

\begin{cousinweb}
\textbf{\sk{पितरौ}} (\emph{pitarau}) is built on \emph{pitṛ}, \textquotedblleft{}father\textquotedblright{} (\textbf{\sk{पिता}}): MW gives its dual \emph{pitarau}, \textquotedblleft{}father and mother, parents\textquotedblright{} (MW).
\end{cousinweb}

\begin{grammarlens}[title={What you've built}]
One more word for describing what you lived and felt.
\end{grammarlens}

\subsection*{Your turn}
\begin{itemize}
\raggedright
  \item \emph{Say it:} \emph{pitarau}
  \item \emph{Say it:} \emph{pitarau}, once more
  \item \emph{Say it:} \emph{ahaṃ pratidinaṃ pitarau namāmi.} — Every day I bow to my parents.
  \item \emph{Recall:} say the Sanskrit for gets along, then the Sanskrit for husband and wife, then say \emph{pitarau} again
\end{itemize}

\begin{tcolorbox}[breakable,colback=teal!4,colframe=teal!35!black,title={Before you move on}]
What does \sk{पितरौ} mean? (\textquotedblleft{}Parents\textquotedblright{}.) Say it once more.
\end{tcolorbox}
\section[(dialogue)]{Looking back: from Growing Up to Getting Along}

\label{lesson:SA-R501-review-495-501}



\begin{quote}
Say the Sanskrit for a lad and youth (the time of life).
\end{quote}

\subsection*{Your turn}
Say these aloud:

\begin{itemize}
\raggedright
  \item a lad, youth (the time of life), grown up, holds back, a celibate student
  \item the marriage thread, a bride, marries, a pavilion, the seven steps
  \item giving away the bride, taking the bride's hand, a fire altar, offers into the fire, good fortune
  \item an elder brother, a younger brother, eldest, youngest, twins
  \item a father's sister, a mother's sister, a brother's son, a sister's son, a daughter's son
  \item a husband's brother, a husband's sister, a wife's brother, a great-grandfather, offspring
  \item a householder, a relationship, gets along, husband and wife, parents
\end{itemize}

\begin{tcolorbox}[breakable,colback=teal!4,colframe=teal!35!black,title={Before you move on}]
A lad? (\textbf{\sk{किशोरः}}, \emph{kiśoraḥ}.) Youth (the time of life)? (\textbf{\sk{यौवनम्}}, \emph{yauvanam}.) Grown up? (\textbf{\sk{प्रौढः}}, \emph{prauḍhaḥ}.) Holds back? (\textbf{\sk{वारयति}}, \emph{vārayati}.) A celibate student? (\textbf{\sk{ब्रह्मचारी}}, \emph{brahmacārī}.) The marriage thread? (\textbf{\sk{मंगलसूत्रम्}}, \emph{maṅgalasūtram}.) A bride? (\textbf{\sk{वधूः}}, \emph{vadhūḥ}.) Marries? (\textbf{\sk{परिणयति}}, \emph{pariṇayati}.) A pavilion? (\textbf{\sk{मण्डपः}}, \emph{maṇḍapaḥ}.) The seven steps? (\textbf{\sk{सप्तपदी}}, \emph{saptapadī}.) Giving away the bride? (\textbf{\sk{कन्यादानम्}}, \emph{kanyādānam}.) Taking the bride's hand? (\textbf{\sk{पाणिग्रहणम्}}, \emph{pāṇigrahaṇam}.) A fire altar? (\textbf{\sk{वेदी}}, \emph{vedī}.) Offers into the fire? (\textbf{\sk{जुहोति}}, \emph{juhoti}.) Good fortune? (\textbf{\sk{सौभाग्यम्}}, \emph{saubhāgyam}.) An elder brother? (\textbf{\sk{अग्रजः}}, \emph{agrajaḥ}.) A younger brother? (\textbf{\sk{अनुजः}}, \emph{anujaḥ}.) Eldest? (\textbf{\sk{ज्येष्ठः}}, \emph{jyeṣṭhaḥ}.) Youngest? (\textbf{\sk{कनिष्ठः}}, \emph{kaniṣṭhaḥ}.) Twins? (\textbf{\sk{यमलौ}}, \emph{yamalau}.) A father's sister? (\textbf{\sk{पितृष्वसा}}, \emph{pitṛṣvasā}.) A mother's sister? (\textbf{\sk{मातृष्वसा}}, \emph{mātṛṣvasā}.) A brother's son? (\textbf{\sk{भ्रातृजः}}, \emph{bhrātṛjaḥ}.) A sister's son? (\textbf{\sk{भागिनेयः}}, \emph{bhāgineyaḥ}.) A daughter's son? (\textbf{\sk{दौहित्रः}}, \emph{dauhitraḥ}.) A husband's brother? (\textbf{\sk{देवरः}}, \emph{devaraḥ}.) A husband's sister? (\textbf{\sk{ननान्दा}}, \emph{nanāndā}.) A wife's brother? (\textbf{\sk{स्यालः}}, \emph{syālaḥ}.) A great-grandfather? (\textbf{\sk{प्रपितामहः}}, \emph{prapitāmahaḥ}.) Offspring? (\textbf{\sk{सन्ततिः}}, \emph{santatiḥ}.) A householder? (\textbf{\sk{गृहस्थः}}, \emph{gṛhasthaḥ}.) A relationship? (\textbf{\sk{सम्बन्धः}}, \emph{sambandhaḥ}.) Gets along? (\textbf{\sk{संगच्छते}}, \emph{saṅgacchate}.) Husband and wife? (\textbf{\sk{दम्पती}}, \emph{dampatī}.) Parents? (\textbf{\sk{पितरौ}}, \emph{pitarau}.)
\end{tcolorbox}
\section[(dialogue)]{Getting Along — once more}

\label{lesson:SA-R501-review-final-501}



\begin{quote}
Say the Sanskrit for a householder and a relationship.
\end{quote}

\subsection*{Your turn}
\begin{itemize}
\raggedright
  \item \emph{Say it:} a householder, a relationship, gets along, husband and wife, parents
\end{itemize}

\begin{tcolorbox}[breakable,colback=teal!4,colframe=teal!35!black,title={Before you move on}]
A householder? (\textbf{\sk{गृहस्थः}}, \emph{gṛhasthaḥ}.) A relationship? (\textbf{\sk{सम्बन्धः}}, \emph{sambandhaḥ}.) Gets along? (\textbf{\sk{संगच्छते}}, \emph{saṅgacchate}.) Husband and wife? (\textbf{\sk{दम्पती}}, \emph{dampatī}.) Parents? (\textbf{\sk{पितरौ}}, \emph{pitarau}.)
\end{tcolorbox}
